\section{The Decision Procedure}
\label{sec:decision}

Section~\ref{sec:costmodel} answers ``how much salt.'' The prior question is
whether to salt at all, given that AQE may already handle the case. The
procedure in Algorithm~\ref{alg:decide} answers both, from statistics available
before execution.

\begin{algorithm}[t]
\caption{Skew mitigation decision procedure. Branch structure and thresholds
are a direct transcription of \texttt{costmodel.decide} and
\texttt{costmodel.recommend\_k} in the released artifact.}
\label{alg:decide}
\begin{algorithmic}[1]
\REQUIRE $H$, $P$, $N$, $M$, $n$, row width $w$, partitions $p$, slots $C$;
         AQE parameters $\beta$ (threshold), $\phi$ (factor), $\alpha$
         (advisory size); optional calibrated $\gamma$
\ENSURE  a strategy, and a salt cardinality where one is required
\STATE $\rho \leftarrow H p / N$
\IF{$\rho < \rho_{\min}$}
    \RETURN \textsc{DoNothing} \COMMENT{no straggler exists}
\ENDIF
\STATE $B_{\text{hot}} \leftarrow H w$;\quad
       $B_{\text{med}} \leftarrow (N-H)\,w/(p-1)$;\quad
       $\Theta \leftarrow \max(\beta,\ \phi B_{\text{med}})$
\IF{$B_{\text{hot}} \le \alpha$}
    \RETURN \textsc{Salt} \COMMENT{no split possible; the rule is a no-op}
\ENDIF
\IF{$B_{\text{hot}} \le \Theta$}
    \RETURN \textsc{Salt} \COMMENT{AQE will not classify the partition as skewed}
\ENDIF
\STATE $\mu \leftarrow P n / M$
\IF{$\mu \le \mu_{\min}$}
    \RETURN \textsc{AqeOnly}
\ELSE
    \RETURN \textsc{Aqe}$+$\textsc{Salt}
\ENDIF
\medskip
\STATE \textbf{if a salt cardinality is required:}
\IF{$\gamma$ is calibrated \textbf{and} the fit is identified}
    \STATE $k^{*} \leftarrow \sqrt{\gamma H / P}$
\ELSE
    \STATE $k^{*} \leftarrow \textsc{Undefined}$ \COMMENT{omit from the minimum}
\ENDIF
\RETURN $k_{\text{eff}} \leftarrow \max\bigl(1, \lfloor\min(k^{*}, \rho, C)\rceil\bigr)$
\end{algorithmic}
\end{algorithm}

The procedure has three branches, and each corresponds to a distinct failure or
success mode.

\subsection{Branch 1: no straggler}
When $\rho < \rho_{\min}$ (we use $\rho_{\min}=2$) the hot key amounts to less
than a couple of ordinary partitions. There is nothing to relieve, and both
mechanisms would add overhead for no gain. This branch exists because
practitioners routinely salt joins that were never skewed enough to matter,
having diagnosed ``skew'' from a slow job rather than from the key
distribution.

\subsection{Branch 2: AQE will not fire}
AQE's rule is conjunctive in \emph{three} conditions, not the two that are
usually discussed. The partition must exceed an absolute byte threshold, a
multiple of the median, \emph{and} the advisory partition size --- because that
last is the size the rule splits into, so a hot partition smaller than it
cannot be divided at all and the optimisation is a no-op however skewed it is.
Section~\ref{sec:results} finds this third condition to be the binding one in
practice, and it is the one practitioners do not tune. A join can be badly skewed in the sense
that matters --- one task running many times longer than its peers --- while
sitting below the absolute threshold, whose default is large. In that regime
AQE does nothing at all, silently, and the practitioner who has enabled it
believes the problem is handled. The remedy is either to salt explicitly or to
lower the threshold; the procedure reports which, rather than leaving the user
to discover the non-event from the Spark UI.

\subsection{Branch 3: the one-sided/two-sided hinge}
This is the substantive branch. When skew is one-sided, splitting the fact-side
partition and replicating the small matching probe-side partition to each split
is close to free, and AQE resolves the straggler. When both sides carry the
same hot key, each of the $s$ splits inherits a large partner. Probe-side work
grows with $s$, and what looked like $s$-fold relief on the critical path is
substantially cancelled.

Detecting which case obtains is cheap, and it is exactly
Equation~\eqref{eq:mu}: compare $P$ against the probe side's mean rows per key.
We use $\mu_{\min}=2$ throughout, in the
text and in the implementation. This is a single group-by on a sample, and it is the
highest-value statistic a practitioner can collect before deciding how to fix a
slow join --- it selects the mechanism \emph{and}, through
Equation~\eqref{eq:kstar}, sets its parameter.

\subsection{Worked example}

Take the configuration measured in Section~\ref{sec:results}: $N=\NFACT$ rows,
$p=\NPART$, $H\approx\HHOT$ rows, $P\approx\PHOT$ rows.

The procedure needs the size of the hot shuffle partition, and here the
analytical estimate and the measurement disagree in a way worth stating.
Estimating $B_{\text{hot}} = Hw$ gives \HOTMB~MB; the event log records the
largest join-stage shuffle partition at \MeasHotMB~MB. The estimate is low
because the hot partition also absorbs whatever cold keys hash to it. In the
other direction, $B_{\text{med}}$ estimates the median at \THRESHMB~MB against
a measured \MeasMedMB~MB, because it ignores variance among the cold
partitions. $Hw$ and $(N-H)w/(p-1)$ are therefore best read as a lower and an
upper bound respectively, and where an event log is available the measured
sizes should be used instead --- as they are in Section~\ref{sec:results}.

Evaluated on the measured sizes, with $\rho=\RHOVAL$ a straggler exists; the
hot partition clears both documented trigger conditions; and it is nonetheless
smaller than the advisory split size, so the procedure returns \textsc{Salt}
with the reason that no split is possible. Section~\ref{sec:results} reports
what Spark actually did.
